\documentclass[11pt]{article}
\usepackage[utf8]{inputenc}
\usepackage{amsmath}
\usepackage{amssymb}
\usepackage{geometry}
\usepackage{booktabs}
\usepackage{hyperref}
\geometry{margin=1in}

\hypersetup{colorlinks=true, linkcolor=blue, urlcolor=blue}

\title{MethylPredictor Theoretical Foundation}
\author{MethylPipeline}
\date{}

\begin{document}
\maketitle
\tableofcontents
\newpage

% -----------------------------------------------------------------------
\section{Goal}

\texttt{MethylPredictor} evaluates a trained classifier (from
\texttt{MethylDetector}/\texttt{MethylClassifier}) on \textbf{holdout test
samples} in two groups: \emph{control} (e.g.\ healthy, class~0) and
\emph{disease} (e.g.\ sick, class~1).  It does not train or modify the model;
it runs classification on the test sets and computes \textbf{classification
metrics}.

% -----------------------------------------------------------------------
\section{Role in the Pipeline}

\begin{description}
  \item[MethylCentroid] Builds cohort centroids per group.
  \item[MethylDetector] Detects DMPs and trains a classifier.
  \item[MethylClassifier] Loads the classifier and predicts on new samples.
  \item[MethylPredictor] Runs MethylClassifier on designated test sets,
    compares predictions to known labels, and reports accuracy metrics.
\end{description}

Samples in the test sets must be \textbf{independent} of those used to build
centroids and train the classifier (true holdout), so that reported metrics
reflect generalisation.

% -----------------------------------------------------------------------
\section{Test Set Labels}

\begin{itemize}
  \item \texttt{test\_control\_paths}: samples assigned expected class~$0$.
  \item \texttt{test\_disease\_paths}: samples assigned expected class~$1$.
\end{itemize}

% -----------------------------------------------------------------------
\section{Classification Metrics}

All metrics use standard scikit-learn definitions and are written to
\texttt{validation\_metrics.json}.

\paragraph{Confusion Matrix.}

For two classes (control~$= 0$, disease~$= 1$):

\[
\begin{pmatrix} TP & FP \\ FN & TN \end{pmatrix}
\]

where $TP$ = true positives (disease predicted as disease), $TN$ = true
negatives (control predicted as control), $FP$ = false positives, $FN$ =
false negatives.

\paragraph{Accuracy.}

\begin{equation}
  \text{Accuracy} = \frac{TP + TN}{TP + TN + FP + FN}.
\end{equation}

\paragraph{Sensitivity (Recall for class~1).}

\begin{equation}
  \text{Sensitivity} = \frac{TP}{TP + FN}.
\end{equation}

\paragraph{Specificity (Recall for class~0).}

\begin{equation}
  \text{Specificity} = \frac{TN}{TN + FP}.
\end{equation}

\paragraph{Balanced Accuracy.}

Appropriate for imbalanced test sets:

\begin{equation}
  \text{Balanced Accuracy} = \frac{\text{Sensitivity} + \text{Specificity}}{2}.
  \label{eq:balanced_acc}
\end{equation}

\paragraph{Precision and F1 Score.}

\begin{equation}
  \text{Precision} = \frac{TP}{TP + FP}, \qquad
  F_1 = \frac{2 \cdot \text{Precision} \cdot \text{Sensitivity}}
             {\text{Precision} + \text{Sensitivity}}.
\end{equation}

\paragraph{Macro and Weighted Averages.}

For multiclass output, macro-averaged precision, recall, and $F_1$ are the
unweighted means over classes; weighted averages weight each class by its
support (number of samples).

% -----------------------------------------------------------------------
\section{Summary}

\begin{table}[h]
\centering
\begin{tabular}{@{}ll@{}}
\toprule
Aspect & Role \\
\midrule
Input & Trained classifier, \texttt{test\_control\_paths},
        \texttt{test\_disease\_paths} \\
Labels & Control $\to 0$, disease $\to 1$ (fixed by group membership) \\
Key metric & Balanced Accuracy, Eq.~\eqref{eq:balanced_acc} \\
Outputs & \texttt{validation\_metrics.json}, \texttt{predictions.csv} \\
Theory & Standard supervised evaluation; no new probabilistic model \\
\bottomrule
\end{tabular}
\caption{Summary of the MethylPredictor evaluation.}
\end{table}

\end{document}
